% Section 3 draft: Replica (testbed). Paste over the current \section{Replica} block.
% EDIT: marks facts to fill in or check before submission.

\section{Replica}
\vspace{0.35em}

The replica is a laboratory version of the onboard detection pipeline. It runs
the same processing chain as the intended system, from video input to annotated
detections, on the edge hardware that is considered for the aircraft. Recorded
video replaces the live camera, and a host computer replaces the operator
console. Both the input and the model can be exchanged without changing the
pipeline, so the replica serves as a test bench for the detection models and
for every optimization in Section~4.

\subsection{Laboratory Platform}

The edge platform is an NVIDIA Jetson Orin Nano Super Developer Kit. It
combines a six-core Arm Cortex-A78AE CPU and an Ampere-architecture GPU with
1024 CUDA cores and 32 Tensor Cores, on one module with 8~GB of 128-bit LPDDR5
memory (102~GB/s)~\cite{nvidia-orin-nano-super}. The GPU is integrated: CPU
and GPU use the same physical memory, so data passed between them does not
cross a PCIe bus, but the two compete for the same memory bandwidth.
% EDIT: cite NVIDIA's CUDA for Tegra application note for the shared memory (not on the product page).
Power modes (\texttt{nvpmodel}) set the number of active CPU cores and the
maximum CPU, GPU and memory (EMC) frequencies. The Super configuration offers
three: 15~W, 25~W (the default) and MAXN SUPER~\cite{nvidia-jetson-power}
(Table~\ref{tab:power-modes}). NVIDIA describes the MAXN modes as
unconstrained but experimental: hardware throttling engages when the module's
power exceeds its budget.

\begin{table}[H]
  \centering
  \small
  \begin{tabular}{@{}lrrrr@{}}
    \toprule
    Mode & CPU cores & CPU max & GPU max & EMC max \\
    \midrule
    15\,W       & 6 & 1498 & 612  & 3199 \\
    25\,W       & 6 & 1344 & 918  & 3199 \\
    MAXN SUPER & 6 & 1728 & 1020 & 3199 \\
    \bottomrule
  \end{tabular}
  \vspace{0.5em}
  \caption{Power modes of the Jetson Orin Nano 8\,GB Super; maximum
  frequencies in MHz.}
  \label{tab:power-modes}
\end{table}
\vspace{-0.6em}
The software stack is Ubuntu 24.04 with the Linux for Tegra kernel, CUDA,
TensorRT for model execution, and PyTorch for the pipeline code around it.
% EDIT: JetPack, CUDA, TensorRT and PyTorch versions (dpkg -l | grep -E "nvidia-l4t-core|tensorrt";
% python3 -c "import tensorrt, torch; print(tensorrt.__version__, torch.__version__)").

Figure~\ref{fig:lab-setup} shows the bench. The Jetson is connected to a
monitor for the annotated output and to the host computer over Ethernet. % EDIT: match the photo.

\begin{figure}[H]
    \centering
    \includegraphics[
        width=\columnwidth,
        totalheight=0.28\textheight,
        keepaspectratio
    ]{figures/lab_setup.jpg} % EDIT: photo file name
    \caption{Laboratory replica: Jetson Orin Nano Super, monitor and host
    computer.}
    \label{fig:lab-setup}
\end{figure}

\subsection{System Architecture and Data Flow}

Figure~\ref{fig:architecture-diagram} shows the data flow. On the Jetson, each
frame passes through five steps: decoding of the 1920$\times$1080 video,
preprocessing (color conversion as the model requires, resizing and
letterboxing to the network input), TensorRT inference, postprocessing (confidence filtering,
non-maximum suppression and mapping of the boxes back to the frame), and
annotation for the display. The input is a recorded video file for the
experiments; the same pipeline also reads an RTSP stream, which is how a
camera feed would arrive. The annotated frames can go to the monitor or be
streamed on to another machine. For the throughput experiments the display is
disabled (Section~5).

Profiling reports are written on the Jetson, and
copied to the host computer, where the analysis scripts and the Nsight
graphical interfaces run. The host takes no part in the measured pipeline.

\begin{figure}[H]
    \centering
    \includegraphics[
        width=\columnwidth,
        totalheight=0.32\textheight,
        keepaspectratio
    ]{figures/diagram.png}
    \caption{System architecture of the replica. Video is decoded,
    preprocessed, run through the TensorRT engine, postprocessed and annotated
    on the Jetson; profiling reports go to the host computer for analysis.}
    \label{fig:architecture-diagram}
\end{figure}

\subsection{Profiling Infrastructure}

Two NVIDIA tools give complementary views of the pipeline.

\textbf{Nsight Systems} (\texttt{nsys}) records a timeline of the whole
process: CPU threads, CUDA API calls, memory copies and GPU kernels. It is a
tracing tool: rather than sampling counters, it records an event with start
and end timestamps for every traced call and GPU operation. For CUDA it builds
on CUPTI, NVIDIA's profiling and tracing interface~\cite{nvidia-nsys-guide},
whose activity interface asynchronously records CUDA API calls, kernels and
memory copies~\cite{nvidia-cupti}. Two trace methods exist: a hardware-based
trace, enabled by default on supported GPUs beginning with Blackwell, and the
older software-instrumented trace; on the Orin's
Ampere GPU the software-instrumented method applies.
% EDIT: confirm in a report's Diagnostics Summary, which states the CUDA trace method used.
NVTX is a header-only annotation library: the application marks named ranges,
and Nsight Systems collects them and places them on the same timeline. The events are buffered during the run and saved to
storage; at the end of the collection the intermediate \texttt{.qdstrm} file
is converted into the \texttt{.nsys-rep} report,
which we export to SQLite for our own analysis. The
pipeline code marks each frame's steps with NVTX ranges (\texttt{decode},
\texttt{inference}, \texttt{draw}, \texttt{display}), where
\texttt{inference} covers preprocessing, the TensorRT engine and
postprocessing. From the exported reports we derive, per frame, the GPU
kernel time inside \texttt{inference}, the idle gaps between kernels, and the
CPU work before and after the kernels. All runs
use \texttt{nsys profile --trace=cuda,nvtx --cuda-graph-trace=node}. By
default, Nsight Systems traces a CUDA graph as a whole, so a replayed graph
appears as one item instead of its individual kernels; node-level tracing
shows each kernel, at the price of extra overhead.

\textbf{Nsight Compute} (\texttt{ncu}) examines single kernels. It inserts
its measurement libraries into the application process and collects metrics
from the GPU each time it detects a kernel launch that matches its filters~\cite{nvidia-ncu-guide}. The metrics come from three sources: the GPU's
hardware performance monitors, software patching of the kernel's
instructions, and launch or device attributes. The
GPU can collect only a limited number of counters at a time, so a large
metric set needs several passes over the same kernel. In kernel replay, the
default, Nsight Compute saves the GPU memory the kernel can access before the
first pass and restores what the kernel wrote before each further pass, so
every pass sees the same input. By default it also
flushes all GPU caches before each replay pass; our
runs keep this default, so the reported cache hit rates and durations describe
a kernel starting from cold caches, not from the warm caches it sees inside
the pipeline. Stall reasons come from warp sampling: at a fixed interval of
cycles, a sampler in each streaming multiprocessor picks an active warp and
records its program counter and scheduler state.
Its reports include occupancy, compute and memory
throughput as a percentage of the hardware's peak (``Speed of Light''),
memory-workload statistics such as cache hit rates, and sampled warp stall
reasons. We use it on the most expensive kernels that Nsight Systems
identifies, with the full section set (\texttt{--set full}), for one frame's
launches (\texttt{--launch-skip}, \texttt{--launch-count}). By default
Nsight Compute locks the GPU clocks during profiling; clock control is not
supported on Orin, and NVIDIA recommends \texttt{jetson\_clocks} on Jetson
instead~\cite{nvidia-ncu-cli}. We therefore pass
\texttt{--clock-control none} and pin the clocks with \texttt{jetson\_clocks},
as in the Nsight Systems runs; the per-launch durations agree with Nsight
Systems to within 3--15\,\%.

\textbf{Overhead.} Neither tool is free, and they disturb the pipeline in
different ways. The cost of Nsight Systems' CUDA tracing is incurred per API
call, and node-level graph tracing adds to it. The
cost therefore depends on how often the pipeline calls the CUDA API and how
CPU-bound it is: a pipeline whose frame time is dominated by GPU kernels loses
little, while a pipeline that issues many short API calls per frame can lose a
substantial part of its throughput. Section~6 quantifies this for both models
and every optimization stage. % EDIT: point to the profiling-overhead subsection label.
For this reason the throughput figures in this paper come from separate runs
without a profiler, and the profiled runs are used only to explain where the
time goes. Nsight Compute serializes the kernel launches of the profiled
application and may replay each profiled kernel over several passes, so it changes the timing of the whole pipeline. It
is never used to measure throughput, only to characterize individual
kernels.

\textbf{Controlled conditions.} Every measurement uses the same setup: the
power mode is set and \texttt{jetson\_clocks} pins the CPU, GPU and memory
clocks at the mode's maximum (at 25~W we read back GPU 918~MHz, CPU
1344~MHz and EMC 3199~MHz, the limits in Table~\ref{tab:power-modes}); the graphical desktop is stopped; and the
over-current counters are read before and after each run to detect
throttling. A run script records these conditions with every result. Without
pinned clocks and with the desktop running, identical code varied between 78
and 101~FPS in our early tests, which hid the effect of individual changes.
